\section{Woche 1}

\subsection{ATB-Störmeldung: Neue Formulare in der App (Ticket \#161)}

    \subsubsection{Aufgabenstellung}
    In der mobilen Anwendung sollten neue Formulare für die Meldung von Störungen an der \ac{ATB} implementiert werden.
    Nur Unternehmen, die für die Verwendung dieses ATB-Formulars registriert sind, können es in der App verwenden.
    Die Datenmodelle für ATB im Disturbances-Microservice sind bereits implementiert, 
    sodass die Formulare in der App nur noch an die entsprechenden Stellen eingebunden werden müssen.

    \subsubsection{Lösung}
    Zunächst habe ich mich in den bestehenden Code eingearbeitet und die breits vorhandenen Datenmodelle für ATB-Störmeldungen im Disturbances-Microservice analysiert.

    Die Formulare wurde mit xamarin Forms implementiert, um eine plattformübergreifende Lösung für die mobile Anwendung zu gewährleisten. 
    Dabei wurden die Eingabefelder und Validierungen gemäß den Anforderungen des ATB-Formulars gestaltet.
    Die Formulare wurden in der mobilen Anwendung implementiert und an die entsprechenden Stellen eingebunden.
    
\subsection{ATB-Störmeldung: Bearbeitung und Übersicht im Web-Portal (Ticket \#185)}

    \subsubsection{Aufgabenstellung}
    ...

